\newpage
\pagestyle{fancy}

\section{Ulteriori indicazioni}
Di seguito vengono elencate ulteriori indicazioni necessarie per garantire un livello di progettazione e un grado di sicurezza ottimali:

\begin{enumerate}
\item Si raccomanda l'installazione della parte inferiore delle prese ad un'altezza compresa tra 0,5 e 1,15\,m da terra;
\item Si raccomanda di non creare depositi di materiali infiammabili e combustibili nel medesimo locale della stazione di ricarica;
\item Al fine di prevenire effetti termici pericolosi, l'isolamento del cavo di alimentazione per la ricarica deve essere di tipo resistente all'usura;
\item Il cavo di alimentazione deve essere verificato (esame visivo) prima di ciascun utilizzo;
\item Qualora il cavo di alimentazione fosse dotato di schermatura metallica, la stessa sarà messa a terra;
\item Il veicolo elettrico deve essere omologato secondo la normativa vigente, mantenuto in efficienza e sottoposto alle revisioni di legge.
\end{enumerate}

Ad ultimazione dei lavori dovrà essere consegnata la documentazione richiesta dal Decreto Ministeriale 22/01/2008 n°37, la Dichiarazione di Conformità (DiCo), integrata con gli allegati obbligatori ed il manuale di manutenzione dell’impianto elettrico (D.Lgs. 81/08)

I quadri elettrici saranno sottoposti alle presenti prove di collaudo:

\begin{itemize}
\item Controllo visivo delle apparecchiature e relativo cablaggio;
\item Verifica della rispondenza agli schemi circuitali e ai dati tecnici;
\item Verifica della presenza della targa di identificazione;
\item Verifica di funzionamento;
\item Verifica dell'efficienza del circuito di protezione;
\item Prova di isolamento dei circuiti principali (di potenza).
\end{itemize}

Gli impianti saranno sottoposti alle seguenti prove e verifiche, come quanto esplicitato dalla norma CEI 64-8/6:

\begin{itemize}
\item Esame a vista che accerti l'integrità dell'impianto e che i componenti siano conformi alle prescrizioni di sicurezza, scelti in modo corretto ed installati secondo normativa, in modo da non compromettere il grado di sicurezza;
\item Prova di continuità dei conduttori di protezione, compresi anche quelli di eventuali collegamenti equipotenziali principale e supplementare;
\item Misura della resistenza dell'isolamento dell'impianto;
\item Verifica dell'efficienza degli interruttori differenziali tramite prova di intervento effettuata con apparecchiatura idonea.
\end{itemize}

\textbf{Di ciascuna prova dovrà essere fornito report con i risultati ottenuti.}

Tutti i materiali dovranno avere caratteristiche atossiche, essere inodori, non igroscopici e privi di amianto e/o qualsiasi altro componente inquinante e non ammesso dalla vigente legislazione.
Al termine dei lavori dovrà essere fornita la seguente documentazione:
\begin{itemize}
    \item Certificazione relativa ai test di resistenza al fuoco rilasciato da laboratorio autorizzato secondo DM 16/02/2007;
    \item Bolla (o documento di trasporto) di consegna del materiale;
    \item Dichiarazione di conformità del produttore nella quale si certifica che il materiale fornito alla ditta installatrice è conforme alle caratteristiche descritte negli elaborati del certificato di prova;
    \item Dichiarazione di corretta messa in opera fornita dalla ditta installatrice.
\end{itemize}


\subsection{Linee guida condivise}
Le presenti linee guida sono state predisposte al fine di fornire un insieme di indicazioni tecniche e organizzative da seguire scrupolosamente per garantire la sicurezza dell’impianto elettrico in oggetto. Esse rappresentano un riferimento fondamentale per assicurare la conformità alle normative vigenti e per prevenire situazioni di rischio, salvaguardando l’incolumità delle persone e l’integrità delle apparecchiature coinvolte.

\subsubsection{Sicurezza delle installazioni e dei relativi dispositivi di protezione
}
Le installazioni e i relativi dispositivi di protezione saranno realizzati a regola d’arte seguendo le norme CEI vigenti al momento della realizzazione dell’impianto stesso.

\subsubsection{Protezioni elettriche}
Gli impianti elettrici essenziali per l’alimentazione dei punti di ricarica saranno realizzati secondo la regola dell’arte. Essi saranno dotati di dispositivi di protezione contro il sovraccarico ed il cortocircuito che consentano un’apertura automatica del circuito di alimentazione.

\subsubsection{Esercizio e manutenzione }
L’esercizio e la manutenzione degli impianti elettrici saranno effettuati secondo quando indicato dalla normativa tecnica applicabile, nei manuali d’uso e manutenzione forniti dai costruttori e dei relativi dispositivi di protezione, ovvero secondo quanto previsto nel piano dei controlli e delle manutenzioni degli impianti. \\
Le operazioni di controllo periodico e gli interventi di manutenzione degli impianti elettrici saranno svolti da personale specializzato al fine di garantire il corretto e sicuro funzionamento. \\
Le operazioni di controllo periodico e gli interventi di manutenzione degli impianti elettrici saranno documentati ed eventualmente messi a disposizione, su richiesta, al competente comando provinciale dei Vigili del Fuoco.

\subsubsection{Messa in sicurezza}
In caso di incendio, al fine di consentire ai soccorritori di intervenire in sicurezza, sarà effettuato il sezionamento di emergenza, in accordo alla normativa tecnica applicabile. \\
Il punto di ricarica sarà dotato di un dispositivo di comando di sgancio di emergenza, costituto da una bobina a minima tensione alimentata a valle dell’interruttore dell’autorimessa soggetta a sgancio, posizionato in una zona segnalata e facilmente accessibile anche agli operatori di soccorso, che sarà implementato nel circuito di sgancio dell’impianto elettrico dell’autorimessa condominiale esistente. \\
Il suddetto dispositivo di comando di sgancio di emergenza deve necessariamente garantire in completa sicurezza il sezionamento dell’impianto elettrico nei confronti delle sorgenti di alimentazione. Qualora tale condizione non fosse rispettata, sarà obbligatorio provvedere all'adeguamento dell'impianto elettrico in questione; fino al completamento dei suddetti interventi, è vietato utilizzare i punti di ricarica per le autovetture elettriche, declinando sin d'ora ogni responsabilità per danni a persone o cose derivanti da un utilizzo improprio o non autorizzato degli stessi.

\subsubsection{Segnaletica di sicurezza}
L’area in cui sono ubicati gli impianti di ricarica sarà segnalata con apposita cartellonistica conforme alla normativa vigente ed alla normativa in materia di sicurezza e salute sui luoghi di lavoro.

\subsubsection{Organizzazione e gestione della sicurezza antincendio 
}
Al fine di ridurre l'insorgere di incendi e la loro propagazione, saranno adottate una serie di misure preventive e protettive:

\begin{itemize}
\item Gli impianti elettrici saranno realizzati a regola d'arte, con materiali autoestinguenti e non propaganti la fiamma;
\item Sarà eseguita la messa a terra di impianti, strutture e masse metalliche, al fine di evitare la formazione di cariche elettrostatiche;
\item Saranno adottati dispositivi di sicurezza quali estintori, in modo tale da garantire lo spegnimento di un incendio che potrebbe essere causato dalle batterie delle autovetture (solitamente con tecnologia a Litio, resta da verificare il tipo di alimentazione delle autovetture che verranno considerate nell'autorimessa oggetto del presente intervento);
\item Sarà garantito il rispetto dell'ordine e della pulizia ove sarà collocato il punto di ricarica e nei vari locali tecnologici;
\item Saranno garantiti controlli sulle misure di sicurezza;
\item Sarà garantita un'adeguata informazione e formazione del personale addetto alla manutenzione.
\end{itemize}

Inoltre, per prevenire gli incendi:

\begin{itemize}
\item Non sarà previsto il deposito e l'utilizzo di materiali infiammabili e facilmente combustibili;
\item Non sarà previsto l'utilizzo di fonti di calore;
\item Non sarà previsto l'utilizzo di fiamme libere ed in tutta l'area sarà vietato fumare.
\end{itemize}

